% =====================================================================
%  Pacotes e layout do documento.
% =====================================================================

% --- Codificacao e idioma --------------------------------------------
\usepackage[T1]{fontenc}
\usepackage[utf8]{inputenc}
\usepackage[brazil]{babel}

% --- Fonte -----------------------------------------------------------
% O padrao oficio pede uma fonte serifada em corpo 12. A Latin Modern
% cumpre isso e existe em qualquer instalacao. Para usar a Times, que
% alguns orgaos exigem, troque a linha abaixo por \usepackage{mathptmx}.
\usepackage{lmodern}

% --- Margens (MRPR: 3cm esquerda, 1,5cm direita) ----------------------
\usepackage[
  a4paper,
  left=3cm,
  right=1.5cm,
  top=2.5cm,
  bottom=2.5cm
]{geometry}

% --- Elementos graficos ----------------------------------------------
\usepackage{graphicx}
\graphicspath{{figures/}}

% --- Espacamento -----------------------------------------------------
\usepackage{setspace}
\usepackage{parskip}
\usepackage{indentfirst}

% --- Links -----------------------------------------------------------
\usepackage[hidelinks]{hyperref}
